\chapter{Methodology}
\label{chap:methodology}

This chapter explains the reasoning behind the assistant before its parts are described. The
assistant is a \ac{rag}\index{retrieval-augmented generation} system: for each question it first
looks for the relevant passages in the company's documents and, when a customer has logged in, in
the data of their policy; only then does it write an answer, based on what it found and nothing
else. The sections that follow present the principles behind the design, the technologies chosen
to respect them at zero cost, the sources the knowledge base is built from and the regulatory
constraints that apply. Chapter~\ref{chap:architecture} then shows how the parts work together,
and Chapter~\ref{chap:knowledge-base} describes the data in detail.

\section{Design Principles}
\label{sec:principles}\index{design principles}

Three principles drive every choice. Each one starts from a weakness of language models and turns
it into a rule that does not depend on the model behaving well.

\paragraph{A constrained flow.} Models are not predictable: the same question can lead to
different steps on different days. For this reason the model only interprets, for example by
recognising what the customer is asking, while plain code decides which step comes next and which
\acsp{api} to call. The path of a question is a fixed graph, and its decisions are rules that can
be read and tested (Section~\ref{sec:routing}).

\paragraph{Identity outside the model.} A prompt can be manipulated: a message can claim to come
from someone else or ask to see another policy. The identity of the customer is therefore never
taken from the conversation. It comes only from the login token (\ac{jwt}), which the policy server
checks on every call, so that a customer can read only their own data whatever the text says
(Section~\ref{sec:components}).

\paragraph{No unverified text.} A draft can be wrong, and a wrong sentence that has already been
shown cannot be taken back. While the answer is being prepared the customer sees only progress
messages; the text appears only after the output guardrail has checked it against its sources
(the \texttt{output\_guard} step in Figure~\ref{fig:graph}).

\section{Technology Selection}
\label{sec:technology-selection}\index{technology selection}\index{Groq}\index{gpt-oss}\index{Qwen3}\index{Gemini}\index{free quota}

Four criteria guided the choice of the technologies. The system must cost nothing: it uses open
source software on one free cloud server, plus the free quotas of external \acsp{api}. Every model
that takes part in writing an answer has open weights under the Apache~2.0 licence; Gemini, a
closed model, is used only to judge the answers during the tests. No \ac{gpu} is needed: the small
models run on the server's \ac{cpu}, and the large ones are reached through an \ac{api}. Finally,
every provider is reached through an OpenAI-compatible \ac{api}, so replacing one is a single line
of configuration.

Figure~\ref{fig:models} compares the models, drawn at the same scale. The two large ones run
through Groq's free \ac{api}: gpt-oss-120b writes the answers, while gpt-oss-20b recognises the
intent of each question and checks the answers. Both are \ac{moe} models, so only a small part of
their parameters is active for each token. The three small Qwen3 models run on the server's
\ac{cpu}, for the search, the reranking of the passages and the screening of the incoming messages.
Mistral \ac{ocr}, finally, is used only to prepare the documents, before the assistant starts.

% Figure 1.2: the models, their role, where they run and their size.
% Bars: 0.085 cm per billion parameters; light = total, dark = active for each token.
\begin{figure}[H]
  \centering
  \begin{tikzpicture}[font=\footnotesize]
    \newcommand\scale{0.085}
    % \model{y}{logo}{name}{role, where it runs}{total B}{active B}{specifications}
    \newcommand\model[7]{%
      \node at (0.35,#1) {\includegraphics[height=0.55cm, width=0.7cm, keepaspectratio]{logos/#2}};
      \node[anchor=west, font=\footnotesize\bfseries] at (0.8,#1+0.17) {#3};
      \node[anchor=west, font=\scriptsize, text=black!60] at (0.8,#1-0.2) {#4};
      \fill[catmodel!22] (5.3,#1+0.05) rectangle ({5.3+#5*\scale},#1+0.3);
      \fill[catmodel] (5.3,#1+0.05) rectangle ({5.3+#6*\scale},#1+0.3);
      \node[anchor=west, font=\scriptsize, text=black!75] at (5.25,#1-0.2) {#7};}
    \node[anchor=west, font=\scriptsize\bfseries, text=black!60] at (0.8,0.75) {Model, role, where};
    \node[anchor=west, font=\scriptsize\bfseries, text=black!60] at (5.25,0.75)
      {Parameters, in billions};
    \draw[black!30] (0,0.5) -- (15.75,0.5);

    \model{0}{openai}{gpt-oss-120b}{answer \textbullet\ Groq~\cite{openai2025gptoss}}{117}{5.1}
      {117B total \textbullet\ 5.1B active}
    \model{-1}{openai}{gpt-oss-20b}{intent, check \textbullet\ Groq}{21}{3.6}
      {21B total \textbullet\ 3.6B active}
    \model{-2}{qwen}{Qwen3-Embedding-0.6B}{embedding \textbullet\ server~\cite{qwen2025embedding}}
      {0.6}{0.6}{0.6B \textbullet\ 32K context \textbullet\ $\leq$1024 dims}
    \model{-3}{qwen}{Qwen3-Reranker-0.6B}{reranking \textbullet\ server~\cite{qwen2025embedding}}
      {0.6}{0.6}{0.6B \textbullet\ 32K context}
    \model{-4}{qwen}{Qwen3Guard-Gen-0.6B}{input guardrail \textbullet\ server~\cite{zhao2025qwen3guard}}
      {0.6}{0.6}{0.6B \textbullet\ 119 languages}
    % Mistral OCR: no logo in logos/, an icon instead; closed model, so no parameter bar
    \node at (0.35,-5) {{\color{catmodel}\Large\faIcon{file-pdf}}};
    \node[anchor=west, font=\footnotesize\bfseries] at (0.8,-5+0.17) {Mistral OCR};
    \node[anchor=west, font=\scriptsize, text=black!60] at (0.8,-5-0.2)
      {reading the PDFs \textbullet\ API~\cite{mistral2026ocr}};
    \node[anchor=west, font=\scriptsize, text=black!75] at (5.25,-5-0.2)
      {closed model \textbullet\ used once, offline};
    \draw[black!30] (0,-5.5) -- (15.75,-5.5);
    % Legend
    \fill[catmodel] (5.3,-5.85) rectangle (5.6,-5.65);
    \node[anchor=west, font=\scriptsize] at (5.65,-5.75) {active for each token};
    \fill[catmodel!22] (8.6,-5.85) rectangle (8.9,-5.65);
    \node[anchor=west, font=\scriptsize] at (8.95,-5.75) {total, for mixture-of-experts models};
  \end{tikzpicture}
  \caption{The models, drawn at the same scale: the large ones run through Groq's free API, the
    small ones on the server. Every model in the answer path is open-weight (Apache~2.0).}
  \label{fig:models}
\end{figure}


The documents were prepared with two more tools, outside the answer path. Mistral
\ac{ocr}~\cite{mistral2026ocr} converted the \acsp{pdf} into structured text once, offline, on
Mistral's free plan. Claude~\cite{anthropic2026claude} was used to anonymize
them, following the outline set by the team (Section~\ref{sec:anonymization}).

\section{Knowledge Base Construction}
\label{sec:kb-summary}\index{knowledge base}

The assistant draws on five kinds of public documents and, for customers only, on their own data
(Figure~\ref{fig:sources}). How the public documents are prepared is described in
Chapter~\ref{chap:knowledge-base}.

% Figure 1.4: the sources of the answers, public and personal
\begin{figure}[H]
  \centering
  \begin{tikzpicture}[font=\footnotesize, panel/.style={rounded corners=5pt, inner sep=5pt}]
    % \doc{row}{icon}{name}{what it holds}
    \newcommand\doc[4]{%
      \node at (0.45,-0.55-#1*0.52) {{\color{catdata}\large #2}};
      \node[anchor=west, font=\footnotesize\bfseries] at (0.85,-0.55-#1*0.52) {#3};
      \node[anchor=west, text=black!75] at (3.35,-0.55-#1*0.52) {#4};}
    \node[anchor=west, font=\footnotesize\bfseries] at (0.15,0) {Public documents
      \textmd{\color{black!60}\ \textbullet\ indexed \textbullet\ for everyone}};
    \doc{0}{\faQuestionCircle}{FAQ}{short answers, also from the corrections}
    \doc{1}{\faIcon{file-signature}}{Forms}{what each one is for and how to send it}
    \doc{2}{\faFileContract}{\acs{dip}}{the pre-contractual summary of each product}
    \doc{3}{\faBook}{\acs{cga}}{the terms and conditions, edition by edition}
    \doc{4}{\faLifeRing}{Support pages}{procedures, deadlines, contacts}
    \begin{scope}[on background layer]
      \draw[panel, draw=catdata!60, fill=catdata!5] (0,0.32) rectangle (10.1,-2.95);
    \end{scope}
    \node[panel, cat sec, text width=5.0cm, align=left, anchor=north west] at (10.45,0.32)
      {{\footnotesize\bfseries Customer data \textmd{\color{black!60}\textbullet\ after login}}\\[5pt]
       \logo[0.9em]{mcp}\enspace Read through the policy server: guarantees, limits, deductibles,
       claims, instalments.\\[4pt]
       {\color{bad}\faBan}\enspace Never indexed, never seen by a guest.};
    % Who reads what
    \node[anchor=west] at (0.15,-3.4) {\faUser\enspace\textbf{Guest}: public documents only};
    \node[anchor=west] at (10.5,-3.4) {\faUserLock\enspace\textbf{Customer}: both};
  \end{tikzpicture}
  \caption{The sources of the answers. A guest gets only general answers; a customer who logs in
    also gets answers on their own policy.}
  \label{fig:sources}
\end{figure}


\section{Regulatory Constraints}
\label{sec:constraints}\index{regulatory constraints}\index{AI Act}\index{personal data}\index{audit}

Under the European \ac{ai} Act~\cite{euaiact2024} the main obligation of this assistant is
transparency: the customer must always know that they are talking to an \ac{ai}. The
assistant should not fall among the high-risk systems, because it does not price products or assess
risks in life and health insurance; this reading still has to be confirmed by the legal office.

Personal data is limited by design. A guest shares none. A customer shares only their own data, and
only after login. Presidio masks the personal identifiers in the text, and only what the answer
needs is sent to the external models; the conversation checkpoints are encrypted
with \ac{aes}, and expired chats are deleted. The knowledge base adds nothing to this: being
synthetic, it contains no real person.

Every answer is recorded together with its sources and the versions of the models that produced
it, in a table where each row is chained to the previous one by a hash, so that any later change
is detected. If the record cannot be written, the answer is not given.
